\chapter{Implémentation et Fonctionnalités}

\section{Workflow Caissier}
Le module "Point de Vente" (POS) est conçu pour guider le caissier à travers un processus rigoureux, minimisant les risques d'erreurs humaines.

\subsection{Ouverture de Session}
Avant de pouvoir enregistrer la moindre transaction, le caissier doit impérativement ouvrir une session.
\begin{enumerate}
    \item \textbf{Identification :} Le caissier se connecte avec ses identifiants personnels.
    \item \textbf{Comptage Initial :} Il doit saisir le montant exact présent en caisse (le fond de caisse).
    \item \textbf{Validation :} Le système crée une nouvelle \textit{Session} en base de données avec le statut \texttt{OPEN}.
\end{enumerate}

\subsection{Opérations Quotidiennes}
Une fois la session ouverte, le caissier accède à l'interface de vente. Les fonctionnalités clés incluent :
\begin{itemize}
    \item \textbf{Encaissement :} Enregistrement rapide des ventes avec calcul automatique du rendu de monnaie.
    \item \textbf{Mouvements de Fond :} Possibilité d'ajouter (apport) ou de retirer (prélèvement) de l'argent pour des raisons opérationnelles, avec justification obligatoire.
    \item \textbf{Visualisation :} Un récapitulatif en temps réel des transactions de la session courante.
\end{itemize}

\subsection{Clôture et Réconciliation}
La fin de journée est l'étape la plus critique :
\begin{enumerate}
    \item \textbf{Verrouillage :} Le caissier initie la fermeture. La caisse passe en mode "Verrouillé", bloquant toute nouvelle transaction.
    \item \textbf{Comptage Aveugle :} Le système demande au caissier de saisir le montant final compté, SANS lui afficher le montant théorique attendu. C'est le principe du "Blind Count".
    \item \textbf{Fermeture :} La session est clôturée. Si un écart existe entre le théorique et le réel, il est enregistré et flagué pour l'administrateur.
\end{enumerate}

\section{Supervision Administrateur}
Le panneau d'administration offre une vue d'ensemble ("God Mode") sur l'activité des agences.

\subsection{Tableau de Bord en Temps Réel}
Les administrateurs peuvent voir l'état de chaque caisse (Ouverte/Fermée) et le caissier actuellement connecté. Des alertes visuelles signalent les sessions ouvertes depuis trop longtemps ou présentant des anomalies.

\subsection{Audit et Rapports}
L'historique complet est accessible :
\begin{itemize}
    \item \textbf{Détail des Sessions :} Accès à la liste exhaustive des mouvements d'une session passée.
    \item \textbf{Traçabilité des Écarts :} Identification rapide des caissiers ou des périodes générant des écarts de caisse récurrents.
\end{itemize}
